% Source: https://artofproblemsolving.com/community/c3h3180052_2022_lmt_spring_speed_round__lexington_mathematical_tournament
% Stable identity: aops-topic-3180052-q25
Put $\theta=\angle C$, so $\angle B=2\theta$ and $\angle A=180^\circ-3\theta$. Because $D$ lies on the internal bisector at $B$ and the perpendicular bisector of $AC$, it is the midpoint of the circumcircle arc $AC$ not containing $B$. Thus $CD=AD=225$, and triangle $BCD$ has angles $\theta,2\theta,180^\circ-3\theta$. The sine rule gives $\dfrac{64}{225}=\dfrac{\sin3\theta}{\sin\theta}=4\cos^2\theta-1$, so $\cos\theta=\dfrac{17}{30}$ (here $0<\theta<60^\circ$). Therefore $BD=225\cdot\dfrac{\sin2\theta}{\sin\theta}=450\cos\theta=255$. Finally $I$ lies between $B,D$. We have $\angle IAD=\dfrac{\angle A}{2}+\theta=90^\circ-\dfrac{\theta}{2}$, while $\angle AID=180^\circ-\angle AIB=90^\circ-\dfrac{\theta}{2}$. Hence $ID=AD=225$. Consequently $\boxed{BI=255-225=30}$.
\begin{asy}
usepackage("amsmath");
size(7cm);
pair B=(0,0),C=(64,0),A=(-80.5,17*sqrt(611)/2),D=(144.5,17*sqrt(611)/2),I=(17,sqrt(611));draw(A--B--C--cycle);draw(A--D--C);draw(B--D);draw(A--I,dashed);dot(I);label("$A$",A,NW);label("$B$",B,SW);label("$C$",C,SE);label("$D$",D,NE);label("$I$",I,W);
\end{asy}
